%% kapa: exam=1 round=30 year=2019 period=2 form=A subject=appraisal_law subject_ko=감정평가관계법규 seq=19 number=19
%% source: Q-Net 감정평가사 기출문제 (cst00309) artlSeq=5205205 file=extracted/1008/5205205/2019년 제30회 감정평가사 1차시험 문제지/2019년 제30회 감정평가사 1차시험 문제지/2019년도 제30회 감정평가사 1차 2교시 A형.hwp
%% conversion: kordoc 4.24.51 -> tools/gen1.py (auto); review=proofread-sonnet; flags=none
\begin{question}{19}
부동산 가격공시에 관한 법령상 지가의 공시에 관한 설명으로 옳은 것은?
\begin{choices}
  \item 개별공시지가의 단위면적은 3.3제곱미터로 한다.
  \item 시장ㆍ군수 또는 구청장은 농지보전부담금 부과대상인 토지에 대해서는 개별공시지가를 결정ㆍ공시하지 아니할 수 있다.
  \item 개별공시지가는 토지시장에 지가정보를 제공하고 일반적인 토지거래의 지표가 되며, 감정평가업자가 개별적으로 토지를 감정평가하는 경우에 기준이 된다.
  \item 표준지에 지상권이 설정되어 있을 때에는 그 지상권이 존재하지 아니하는 것으로 보고 표준지공시지가를 평가하여야 한다.
  \item 표준지에 대한 용도지역은 표준지공시지가의 공시사항에 포함되지 않는다.
\end{choices}
\end{question}
\answer{4}
